\itemhead{p.84 练习 · 第3题（教材p.84）}{XJ-B1-C03-PR-P084-Y1488}
\label{XJ-B1-C03-PR-P084-Y1488}
{\footnotesize 湘教版教材p.84 / PDF第91页 / p.84 练习 / 第3题}
\field{原题}{求证：定义于 $\mathbb R$ 上的两个奇函数的乘积是偶函数。}
\field{参考答案}{设 $h(x)=f(x)g(x)$，则 $h(-x)=(-f(x))(-g(x))=h(x)$，定义域为 $\mathbb R$，故为偶函数。}
{\small \field{关联知识}{F08 奇偶性、对称与函数性质应用（教材帮p.113、114、115 / PDF122、123、124）}
\field{方法与易错点}{检查定义域，再代入 $-x$；求另一侧解析式时同时写出 $-x$ 所在区间。 易错：只看公式不看定义域；用一对点证明奇偶；负半轴代换遗漏最外层负号。}
\field{答案依据}{答案册相关题组：PDF第23页（册内页22）；按当前题库原题逐项重算，不按旧版题号机械对套。}}
